\section{Prefix-tree symmetry and hierarchical energy}\label{prefix:section}
The two-sweep estimates restart from every deterministic configuration.
A different gain is available at the end of a completed sweep: the law
of the signed mass has a symmetry that an arbitrary configuration need
not have.  We use this symmetry to contract at each subsequent sweep.
The exact recursion below records the contribution of every coordinate
block, so it also describes which scales retain the energy.

The coordinate order must be the same in successive sweeps. We use
$1,\ldots,d$, with $n=2^d$, and the signed-weight process of
Section~\ref{cycles:weights}. Write $H$ for its available set, $f$
for its signed vector, and $\mathcal E(H,f)=\sum_xf(x)^2$.
The update preserves $|H|$ and $\sum_xf(x)$ and does not increase
$\mathcal E$. Our deterministic recursion refines the common
block-merging identity from Section~\ref{cycles:block-core}. The
probabilistic step will apply that recursion to an invariant random
input law, rather than restart from an arbitrary realized input.

\begin{theorem}[Contraction after the first sweep]\label{prefix:contraction}
Let $H\subseteq\{0,1\}^d$ be deterministic, with $|H|\ge n/4$, and
let $f$ be a deterministic real vector supported on $H$ with
$\sum_x f(x)=0$.  Evolve $(H,f)$ by the conditional signed-mass rule,
using fresh independent switches in the same coordinate order in each
sweep.  If $\mathcal E_b$ is the energy after $b$ sweeps, then, for
all sufficiently large $d$ and every integer $b\ge1$,
\begin{equation}\label{prefix:energy}
 \mathbb E\mathcal E_b\le n^{-(b-1)/64}\mathcal E_0.
\end{equation}
The threshold for $d$ is independent of $H$, $f$ and $b$.
\end{theorem}

The first step is an exact identity, valid without the lower bound on
$|H|$.  We then identify the symmetry of the output and use it to
control the terms in that identity.

\subsection{The exact recursion over coordinate blocks}
A coordinate block of level $j$ leaves the first $j$ bits free and
fixes the remaining bits; its size is $2^j$.  A non-singleton block
$C$ has children $C_0,C_1$ obtained by fixing bit $j$ as well.  For
deterministic initial data $(H,f)$, write
\[
 m_C=|H\cap C|,\qquad \alpha_C=\frac{m_C}{|C|},\qquad
 S_C=\sum_{x\in C}f(x).
\]
Counts and signed sums remain fixed while updates stay inside $C$.
Set $S_C^2/m_C=0$ when $m_C=0$, since then $S_C=0$.  For a
non-singleton block define
\begin{equation}\label{prefix:detail}
 D_C=\frac{S_{C_0}^2}{m_{C_0}}+
     \frac{S_{C_1}^2}{m_{C_1}}-\frac{S_C^2}{m_C}\ge0.
\end{equation}
The inequality is the weighted Cauchy--Schwarz inequality.  The quantity
$D_C$ measures the difference between the two children's average signed
masses per available position.

Let $E_C$ be the expected energy in $C$ after its first $j$ coordinate
updates, and put $W_C=E_C-S_C^2/m_C$.  All the quantities on the
right below refer to the fixed initial data.

\begin{lemma}[Hierarchical energy identity]\label{prefix:hierarchy}
For every non-singleton coordinate block,
\begin{equation}\label{prefix:recursion}
 W_C=(1-\alpha_{C_1}/2)W_{C_0}
     +(1-\alpha_{C_0}/2)W_{C_1}+(1-\alpha_C)D_C.
\end{equation}
If the total signed sum is zero, the expected energy after one sweep,
denoted $\Phi(H,f)$, is therefore
\begin{equation}\label{prefix:expanded}
 \Phi(H,f)=\sum_{C\text{ non-singleton}}(1-\alpha_C)D_C
  \prod_{A\supsetneq C}
    \left(1-\frac{\alpha_{\operatorname{sib}_A(C)}}2\right).
\end{equation}
Here the product runs over the proper coordinate-block ancestors of
$C$, and $\operatorname{sib}_A(C)$ is the child of $A$ not containing
$C$.  Empty products equal one, and
\begin{equation}\label{prefix:telescoping}
 \sum_{C\text{ non-singleton}}D_C=\mathcal E(H,f).
\end{equation}
\end{lemma}
\begin{proof}
Write $a=|C_0|=|C_1|$, and temporarily shorten subscripts to $0,1$.
Since $\alpha_i=m_i/a$, Lemma~\ref{sweeps:merge} reads
\[
 E_C=(1-\alpha_1/2)E_{C_0}
       +(1-\alpha_0/2)E_{C_1}+S_0S_1/a.
\]
We separate from each energy the squared mean on its available sites.
When both counts are positive,
\[
 D_C=\frac{m_0m_1}{m_0+m_1}
         \left(\frac{S_0}{m_0}-\frac{S_1}{m_1}\right)^2,
\]
and direct multiplication gives
\[
 \alpha_C D_C=\frac12\left(
       \alpha_1\frac{S_0^2}{m_0}+\alpha_0\frac{S_1^2}{m_1}\right)
       -\frac{S_0S_1}{a}.
\]
Subtracting the baseline $S_C^2/m_C$ proves
\eqref{prefix:recursion}.  If a child is empty, $D_C=0$ and the
same identity holds using the declared zero convention.

At a singleton $W_C=0$.  Expanding the recursion from the root down
therefore gives \eqref{prefix:expanded}; at the root the baseline
vanishes because the total signed sum is zero.  Finally every
non-root baseline in the sum of \eqref{prefix:detail} occurs once
with each sign.  The surviving leaf terms sum to $\sum_xf(x)^2$,
and the root term is zero.  This proves \eqref{prefix:telescoping}.
\end{proof}

The identity separates energy created by a difference of child means
from the attenuation supplied by the larger blocks above it.  To use
that attenuation repeatedly, we need balance at many scales in the
law entering a sweep.  The following output symmetry supplies it.

\subsection{The symmetry supplied by a completed sweep}
Let $\mathcal A_d$ be the finite group of automorphisms of the binary
prefix tree.  Concretely, its elements have the form
\begin{equation}\label{prefix:automorphism}
 (gx)_i=x_i\oplus F_i(x_1,\ldots,x_{i-1}),\qquad 1\le i\le d,
\end{equation}
where each $F_i$ is an arbitrary function to $\{0,1\}$, including
a constant when $i=1$.  The map is invertible by solving its coordinates
in order.  Write $gf$ for the transported vector,
$(gf)(gx)=f(x)$.

\begin{lemma}[Prefix-tree invariance]\label{prefix:invariance}
For every deterministic initial $(H,f)$, the output law of one sweep
is invariant under $(H',f')\mapsto(gH',gf')$ for every
$g\in\mathcal A_d$.
\end{lemma}
\begin{proof}
Fix the functions in \eqref{prefix:automorphism}. Let $g_{<i}$
transform only coordinates $1,\ldots,i-1$ according to those functions,
leaving coordinates $i,\ldots,d$ unchanged. It bijects the
coordinate-$i$ matching edges. If an old edge $e$ has prefix
$u=(x_1,\ldots,x_{i-1})$ and coin $\omega_i(e)$, define the new
coin by
\[
 \omega_i'(g_{<i}e)=\omega_i(e)\oplus F_i(u).
\]
The prefix here is the original one, recoverable from its transformed
image because $g_{<i}$ is invertible. This is a deterministic
bijection and toggling of the layer's independent fair bits, so all
modified layer arrays have their original joint law.

Inductively the input to the modified layer is the old input with
the earlier coordinates transformed. At an edge with one available input
position, the toggled coin
transports that position and its weight to the image under the first
$i$ transformed coordinates. At an edge with two available input
positions the two new weights are equal, so the same assertion holds.  It is immediate on an empty
edge.  After the last layer the modified output is precisely the
$g$-image of the original output.  Since the coin laws agree, their
output laws agree.
\end{proof}

A uniform element of $\mathcal A_d$ can be generated by independent
fair choices of which two children to interchange at each tree node.
The prefix subtrees and the coordinate blocks cut across each other:
a level-$j$ coordinate block contains exactly one leaf with each
length-$j$ prefix.  This gives the independence in the next lemma.

\begin{lemma}[Independent leaf sampling]\label{prefix:sampling}
Fix deterministic zero-sum data $(H,f)$ with $|H|\ge cn$, where
$c=1/4$, and put $s=2^{\lfloor d/2\rfloor}$.  For uniform
$g\in\mathcal A_d$, with probability at least
$1-2n e^{-cs/8}$, every coordinate block of size at least $s$ has
a fraction at least $c/2$ of its positions in $gH$. For every block $C$ of size
$s$,
\begin{equation}\label{prefix:block-second-moment}
 \mathbb E_g S_C(gf)=0,\qquad
 \mathbb E_g S_C(gf)^2\le(s/n)\mathcal E(H,f).
\end{equation}
\end{lemma}
\begin{proof}
First fix a coordinate block of size $u=2^j$.  Condition on the
automorphism's action through depth $j$.  Each source subtree with
fixed first $j$ bits maps to one target subtree of that depth.  The
actions inside these subtrees are independent uniform automorphisms.
Thus the preimages of the block's $u$ leaves select one uniform
leaf independently from each source subtree, each of size $n/u$.

The hidden count $X$ in the transformed block is a sum of independent
indicators with mean $z=|H|u/n\ge cu$.  Multiplication of their
exponential moments gives
$\mathbb E_g 2^{-X}\le e^{-z/2}$.  Markov's inequality yields
\[
 \Pr_g(X<cu/2)\le\Pr_g(X\le z/2)
 \le 2^{z/2}e^{-z/2}\le e^{-z/8}\le e^{-cu/8}.
\]
There are fewer than $2n$ coordinate blocks, so a union bound proves
the density assertion.  For the signed sum, the sum of the sampled
means is $(u/n)\sum_xf(x)=0$.  Independence gives variance at most
the sum of the sampled second moments, namely
$(u/n)\sum_xf(x)^2$.  Taking $u=s$ proves
\eqref{prefix:block-second-moment}.
\end{proof}

\begin{proof}[Proof of Theorem~\ref{prefix:contraction}]
First randomize the input by uniform $g$ and consider the density event
in Lemma~\ref{prefix:sampling}. On this event,
every term of \eqref{prefix:expanded} with $|C|\le s$ has
$\log_2(n/s)$ ancestors of sizes $2s,4s,\ldots,n$. Their sibling
blocks have size at least $s$, hence density at least $c/2$.  Its multiplier is therefore at most
$(1-c/4)^{\log_2(n/s)}$.  By \eqref{prefix:telescoping}, the sum of
all its nonnegative $D_C$ terms is at most $\mathcal E(H,f)$.

For $|C|>s$, discard the multipliers.  The corresponding details
telescope to $\sum_{|C|=s}S_C^2/m_C$, since the total signed sum is
zero.  On the density event this is at most
$(2/(cs))\sum_{|C|=s}S_C^2$.  There are $n/s$ such blocks; their
second-moment bounds in \eqref{prefix:block-second-moment} sum to
at most $\mathcal E(H,f)$.  On the exceptional event use the
unconditional bound $\Phi\le\mathcal E(H,f)$.  Hence
\begin{equation}\label{prefix:randomized-bound}
 \mathbb E_g\Phi(gH,gf)
 \le\left[(1-c/4)^{\log_2(n/s)}+\frac{2}{cs}
                   +2n e^{-cs/8}\right]\mathcal E(H,f)
 \le n^{-1/64}\mathcal E(H,f)
\end{equation}
for sufficiently large $d$.  Indeed the first term is at most
$e^{-cd/8}$, whose exponent in $n$ is $c/(8\log2)>1/64$,
and $s\ge\sqrt{n/2}$ makes the other terms smaller than that power
eventually.  This threshold depends on no input data.

At each boundary after at least one sweep, the law of $(H,f)$ is
invariant under an independent uniform $g$ by
Lemma~\ref{prefix:invariance}.  The next sweep uses fresh switches
in the same order.  Thus its expected output energy is
\[
 \mathbb E\Phi(H,f)=\mathbb E_{(H,f)}\mathbb E_g\Phi(gH,gf)
                 \le n^{-1/64}\mathbb E\mathcal E(H,f).
\]
This uses invariance of the boundary law, not invariance of each
realized configuration.  The first sweep does not increase energy.
Iteration proves \eqref{prefix:energy} for every $b\ge1$.
\end{proof}

\subsection{Uniform list entropy and the four-block application}
The new estimate gives relative entropy as well as total variation
for every fixed list.  Relative entropy below uses natural logarithms:
$D(p\|q)=\sum_xp(x)\log(p(x)/q(x))$, with $0\log0=0$.

Write $U_{\mathrm{inj},k}$ for the uniform law on ordered injections
of $k$ labels into $V$.

\begin{proposition}[Three-quarter list entropy]\label{prefix:entropy}
For every integer $b\ge1$, every deterministic starting deck, and
any ordered list of $k\le3n/4$ distinct labels, the endpoint law
$\mu_{b,k}$ after $b$ sweeps satisfies, for sufficiently large $d$,
\begin{equation}\label{prefix:entropy-bound}
 D(\mu_{b,k}\|U_{\mathrm{inj},k})
          \le n^2 n^{-(b-1)/64}.
\end{equation}
In particular at $b=769$, every such list has total-variation error
at most $n^{-5}$.
\end{proposition}
\begin{proof}
Let $Y_1,\ldots,Y_k$ be the image positions.  The entropy chain rule
expresses the deficit from $\log(n)_k$, where
$(n)_k=n(n-1)\cdots(n-k+1)$, as
\[
 \sum_{i=1}^k\bigl[\log(n-i+1)-\mathsf H(Y_i\mid Y_1,\ldots,Y_{i-1})\bigr].
\]
Here $\mathsf H$ is Shannon entropy.  Refining the conditioning to
the complete paths of the preceding labels can only lower this
conditional entropy.  Their endpoint complement has
$m=n-i+1\ge n/4$ sites.  For the conditional probability $p$ there,
$\log z\le z-1$ gives
\[
 D(p\|U_H)=\sum_{x\in H}p(x)\log(mp(x))
              \le m\sum_{x\in H}(p(x)-1/m)^2.
\]
The initial centered tag energy is $1-1/m\le1$.  Applying
Theorem~\ref{prefix:contraction} at each exposure and summing at
most $n$ contributions, each with $m\le n$, proves
\eqref{prefix:entropy-bound}.  This also covers the empty list.

For completeness, $D(p\|q)\ge2\|p-q\|_{\mathrm{TV}}^2$ follows
by the log-sum inequality applied to $\{p>q\}$ and its complement.
The resulting binary relative entropy, as a function of its first
mass, has second derivative at least $4$, minimum zero at the
reference mass, and zero first derivative there.  Integrating this
bound, with endpoint values obtained by continuity, proves the
inequality.  With $b=1+12\cdot64=769$, the entropy bound is
$n^{-10}$ and the claimed variation bound follows.
\end{proof}

Let $\mu$ now be the full permutation law after $769$ sweeps.
Partition the labels into four equal blocks, and let $H_i$ permute
block $i$, fixing all other labels.  The left coset $gH_i$ is
specified by the images of the complement of that block.  Each coset
marginal has variation error at most $\epsilon=n^{-5}$ by
Proposition~\ref{prefix:entropy}.  Delete the outcomes belonging
to a coset whose mass exceeds twice its uniform mass, in any of the
four systems.  The remaining subprobability $\mu_0\le\mu$ has
mass $p_0\ge1-8\epsilon$ and satisfies
\begin{equation}\label{prefix:cap}
 \mu_0(gH_i)\le\frac2{[S_n:H_i]}\qquad(i=1,\ldots,4).
\end{equation}
Indeed one system deletes mass at most $2\epsilon$, since the mass
in a heavy coset is at most twice its positive excess over uniform.

We use the four-block type-counting transfer of
\cite[Corollary~\ref*{l-l:four-block-reservoir}]{partialFourier}.
In its precise form needed here, a subprobability satisfying
\eqref{prefix:cap} has, on the irreducible representation
$\rho_\lambda$ of dimension $D_\lambda$,
\begin{equation}\label{prefix:four-block-input}
 \left\|\sum_g\mu_0(g)\rho_\lambda(g)\right\|_{\mathrm{op}}
 \le \sqrt{2R_\lambda}\,D_\lambda^{-1/4}
 \le D_\lambda^{-1/8},
 \qquad \lambda\notin\{(n),(1^n)\},
\end{equation}
for all sufficiently large $n$.  Here $R_\lambda$ is the number of
distinct joint types in the restriction to $H_1\times\cdots\times H_4$,
without counting their multiplicities.  The first inequality retains
one orbit span per joint type.  The second uses
$R_\lambda\le(\sum_{j=0}^{\ell}p(j))^4$ with
$\ell=n-\max(\lambda_1,\lambda'_1)$, and the uniform growth
$\log D_\lambda/\sqrt\ell\to\infty$.
These are the type-counting and degree statements proved in that
companion; an eight-block orbit estimate alone would not give this
four-block bound.

We now check the mass, sign and time substitutions that complete this
application.  Set $q=24$ and
$A_\lambda=\sum_g\mu_0(g)\rho_\lambda(g)$.  Plancherel with
unnormalized trace gives
\[
 |S_n|\sum_g\left|\mu_0^{*q}(g)-\frac{p_0^q}{|S_n|}\right|^2
       =\sum_{\lambda\ne(n)}D_\lambda\|A_\lambda^q\|_{\mathrm{HS}}^2.
\]
For the nontrivial nonsign representations,
\eqref{prefix:four-block-input} bounds the right side by
\[
 \sum_{\lambda\notin\{(n),(1^n)\}}
             D_\lambda^{2-q/4}
   =\sum_{\lambda\notin\{(n),(1^n)\}}D_\lambda^{-4}=o(1),
\]
using the inverse-fourth degree sum from
\cite[Theorem~\ref*{l-l:degree-all-powers}]{partialFourier}.
No normality is required: use
$\|A_\lambda^q\|_{\mathrm{HS}}\le
\sqrt{D_\lambda}\|A_\lambda\|_{\mathrm{op}}^q$.
The original law $\mu$ has zero sign expectation because toggling
one fixed fair switch reverses its permutation parity.  Consequently
$|A_{(1^n)}|\le1-p_0$, and the sign term tends to zero too.
Cauchy--Schwarz now shows that $\mu_0^{*q}$ is $o(1)$ away in
half-$\ell^1$ from $p_0^q U_{S_n}$.

Finally $\mu^{*q}-\mu_0^{*q}$ is nonnegative, of mass $1-p_0^q$.
Its half-$\ell^1$ norm and that of $(1-p_0^q)U_{S_n}$ sum to
$1-p_0^q=o(1)$.  Thus $\mu^{*24}$ converges to uniform in total
variation.  Each factor consists of $769$ complete sweeps, so the
physical time is
\[
 24\cdot769d=18456d.
\]
All statements are uniform over the deterministic starting deck.
Together with the support lower bound \eqref{mart:lower-exact}, this gives
$2d-O(1)\le t_{\mathrm{mix}}(d)\le18456d$ for sufficiently large $d$.
